\section{Complexifying demand-only shocks}\label{sec:beyond-demand-only}

\subsection{Kaldor--Verdoorn effects}\label{subsec:pricing-capacity}

A key observation from the preceding section is that the one-factor model shows limited price responses, which appear only where the labor market is sector-specific: in Table~\ref{tab:results-matrix} the average real wage $\omega$ moves only in the no-reallocation and elastic-supply columns, where sectoral wages respond to the program, and stays at its baseline in the other nine cells. The nature and scope of a transformational stimulus --- especially in a long-term setting --- nevertheless invite reflection on how a permanent shift in demand could affect production in the short and medium run. Additional demand may let firms develop novel capabilities and exploit economies of scale, or it may induce capacity pressures that reduce efficiency: either way the effect works through technology rather than through the demand vector itself.

The most straightforward way to express this intuition formally is to let
effective technology respond to utilization,
\begin{equation}\label{eq:capacity-extension}
  \mathcal{A}_i^{H} \;=\; A_i^{H}\,\Bigl(\frac{y_i}{\lambda_i}\Bigr)^{-\delta\,\alpha_i},
\end{equation}
with $\lambda_i$ as the sector's baseline gross output --- in the baseline-GDP units of the calibration (Appendix~\ref{app:pipeline-va}) the Domar weight, so that $y_i/\lambda_i = 1$ at the baseline --- and $\alpha_i$ as the value-added share, so that $\delta = 0$ reproduces the model of the preceding sections exactly.

Here $y_i$ denotes aggregate sector output. Its effect on productivity is external to individual competitive firms, which take $\mathcal{A}_i^{H}$ as given when choosing inputs and output. Constant returns therefore hold conditional on sectoral productivity, and competitive pricing remains $p_i=c_i$. The elasticity $\delta$ is a reduced-form utilization elasticity of \emph{value-added} productivity: no factor is added, because the productivity change passes into unit cost under zero profit, weighted by the factor share. A Hicks-neutral gain $d\ln \mathcal{A}_i^{H}$ then raises value added per worker by $d\ln \mathcal{A}_i^{H}/\alpha_i$, so $-\delta$ is directly comparable with an estimated Verdoorn coefficient.\footnote{Applied to gross output instead, the same bracket would overstate the value-added response by $1/\alpha_i$, which matters where factor shares are small ($\alpha_i \approx 0.04$ in coke and refined petroleum).}

The two signs of $\delta$ again carry two traditions, underpinned by two distinct visions of the typical operating mode of an economy. With $\delta > 0$ scarcity gains an additional bite: unit costs rise with sector output --- capacity pressure, tight supply chains, higher risk and, hence, eventually higher costs and prices (even if, as here, market power considerations are bracketed out). With $\delta < 0$, on the other hand, unit costs fall with output --- the dynamic increasing returns articulated early by \textcite{Verdoorn1949} (see also \citealp{Kaldor1966} and \citealp{SylosLabini1995}).\footnote{With $\delta < 0$ positive feedback from decreasing costs can give rise to multiple equilibria. We therefore compute all solutions with $\delta \neq 0$ by numerical continuation from $\delta = 0$ and recompute each from four different initial values. In all cases reported, the four solutions agree to within $10^{-6}$ and the Jacobian is non-singular, so the equilibrium is considered locally unique (see Appendix~\ref{app:pricing-hooks}).} Since the model is static, $\delta$ is a level elasticity; Verdoorn's law proper relates growth rates, and~\eqref{eq:capacity-extension} is its static analogue.

We run scenarios with $\delta = \pm 0.5$ in the goods sectors to illustrate the impact of the Verdoorn coefficients on estimated outcomes. Application of these coefficients is confined to the goods sectors, following \textcite{Kaldor1966}, who reads Verdoorn's law as a law of manufacturing, not of services. This setup also matches the incidence of the program, whose direct demand falls entirely on goods sectors. Appendix~\ref{app:profiles} reports more details and a counterfactual that extends the elasticity to services.

Table~\ref{tab:pricing-kaldor} reports the results in the same format as Section~\ref{sec:application}; it covers the full matrix of approaches discussed for scenarios with endogenously increasing ($\delta = -0.5$) and decreasing ($\delta = +0.5$) productivity effects. A first observation is that the employment adjustment model and the Leontief closure no longer coincide, so that we now arrive at twelve distinct economies. Moreover, the average real wage now responds in the wage-pinned panels as well: it falls under capacity pressure ($\delta > 0$) and rises under increasing returns ($\delta < 0$), by about $0.3\%$ in either direction under external finance; the standard CGE and elastic-supply panels move in the same direction. The no-reallocation model behaves differently: there $\omega$ averages sectoral wages, which move with the prices of the sectors that pay them, so it rises slightly under capacity pressure and falls slightly under increasing returns. These changes remain marginal; nonetheless the most pronounced \emph{sectoral} price changes, captured by $\max_i |p_i/\mathcal{P} - 1|$, are also concentrated in the no-reallocation and elastic-supply panels. 

These price changes reflect efficiency gains and losses due to the additional demand, with different implications for production and employment across different model specifications. For labor market closures where scarcity governs economic reactions -- that is, the first three closure specifications in our matrix -- increasing technological capabilities ease these scarcity constraints so that production can expand. For the output adjustment closures with a pinned money wage on the other hand,  increasing efficiency leads to lower employment effects and, correspondingly, (slightly) lower impacts on net welfare.

The different financial closures have close to no impact on how the economy responds to the introduction of Kaldor--Verdoorn effects; in general results reproduce the overall properties identified in Section~\ref{sec:application}. The overall impact of the program is again more pronounced for the quantity adjustment models with a pinned wage. In the no-reallocation, standard-CGE and elastic-supply closures, tax finance (F2) and external finance (F3) continue to produce identical real outcomes when productivity responds to output. As noted above, the perfect correspondence between the employment-adjustment and Leontief columns breaks for any $\delta \neq 0$ as substitution elasticities matter again and prices move. In general, the employment-adjustment model thereby shows greater responsiveness to technological change as compared to the Leontief closure: under external finance increasing returns ($\delta = -0.5$) reduce the employment effect from its $1.78\%$ baseline to $1.35\%$ under employment adjustment, against $1.48\%$ under Leontief. 

% =====================================================================
% Table 6 --- the capacity extension by labor closure and financing (GENERATED)
% ---------------------------------------------------------------------
% Do NOT edit by hand. Regenerate with:  make_pricing_door.py
% OpenCGE tag v0.5.0 (ADR-0044); source run ids are listed below.
% Sources: pricing_capacity_rows-<row>-<fin>-<variant>,
% pricing_capacity-GAMMA-<fin>-<variant>, pricing_mobile-ALPHA-<fin>-<variant>.
% Goods-only profile (the main-text calibration), all five labor rows
% (ADR-0032). The all-sector counterfactual lives in tableA8_pricing-profiles.tex.
% COLLAPSED CELLS: in the three flexible-wage panels (no-reallocation,
% standard CGE, elastic supply) tax finance (F2) and external finance (F3)
% coincide, so their entries are printed once, in one cell spanning both
% groups; the two wage-pinned panels print them separately. The generator
% asserts the coincidence on the PRINTED strings, so a broken equivalence
% aborts the build.
% =====================================================================
\begin{table}[htbp]
\centering
\footnotesize
\definecolor{tableguide}{gray}{0.92}
\definecolor{tableline}{gray}{0.72}
\setlength{\tabcolsep}{3pt}
\renewcommand{\arraystretch}{1.2}
\resizebox{\textwidth}{!}{%
\ifdefined\pdwA\else\newlength{\pdwA}\fi
\settowidth{\pdwA}{$+0.62$}\settowidth{\dimen0}{$\Delta\omega$}\ifdim\dimen0>\pdwA\setlength{\pdwA}{\dimen0}\fi
\ifdefined\pdwB\else\newlength{\pdwB}\fi
\settowidth{\pdwB}{$22.14$}\settowidth{\dimen0}{\shortstack{$\max$\\$|\Delta\hat p|$}}\ifdim\dimen0>\pdwB\setlength{\pdwB}{\dimen0}\fi
\ifdefined\pdwC\else\newlength{\pdwC}\fi
\settowidth{\pdwC}{$+0.12$}\settowidth{\dimen0}{$\Delta L$}\ifdim\dimen0>\pdwC\setlength{\pdwC}{\dimen0}\fi
\ifdefined\pdwD\else\newlength{\pdwD}\fi
\settowidth{\pdwD}{$-0.01$}\settowidth{\dimen0}{$\Delta Y$}\ifdim\dimen0>\pdwD\setlength{\pdwD}{\dimen0}\fi
\begin{tabular}{l@{\hspace{3pt}}p{2pt}@{\hspace{3pt}}>{\raggedleft\arraybackslash}p{\pdwA}!{\color{tableline}\vrule width 0.3pt}>{\raggedleft\arraybackslash}p{\pdwB}!{\color{tableline}\vrule width 0.3pt}>{\raggedleft\arraybackslash}p{\pdwC}!{\color{tableline}\vrule width 0.3pt}>{\raggedleft\arraybackslash}p{\pdwD}@{\hspace{3pt}}p{2pt}@{\hspace{3pt}}>{\raggedleft\arraybackslash}p{\pdwA}!{\color{tableline}\vrule width 0.3pt}>{\raggedleft\arraybackslash}p{\pdwB}!{\color{tableline}\vrule width 0.3pt}>{\raggedleft\arraybackslash}p{\pdwC}!{\color{tableline}\vrule width 0.3pt}>{\raggedleft\arraybackslash}p{\pdwD}@{\hspace{3pt}}p{2pt}@{\hspace{3pt}}>{\raggedleft\arraybackslash}p{\pdwA}!{\color{tableline}\vrule width 0.3pt}>{\raggedleft\arraybackslash}p{\pdwB}!{\color{tableline}\vrule width 0.3pt}>{\raggedleft\arraybackslash}p{\pdwC}!{\color{tableline}\vrule width 0.3pt}>{\raggedleft\arraybackslash}p{\pdwD}}
\toprule
 & & \multicolumn{4}{c}{\textbf{F1}} & & \multicolumn{4}{c}{\textbf{F2}} & & \multicolumn{4}{c}{\textbf{F3}} \\
 & & \multicolumn{4}{c}{\scriptsize repurposing of consumption} & & \multicolumn{4}{c}{\scriptsize refinancing via taxes} & & \multicolumn{4}{c}{\scriptsize external finance} \\
\cmidrule(lr){3-6}\cmidrule(lr){8-11}\cmidrule(lr){13-16}
\multicolumn{1}{c}{$\delta$} & & \multicolumn{1}{c}{$\Delta\omega$} & \multicolumn{1}{c}{\shortstack{$\max$\\$|\Delta\hat p|$}} & \multicolumn{1}{c}{$\Delta L$} & \multicolumn{1}{c}{$\Delta Y$} & & \multicolumn{1}{c}{$\Delta\omega$} & \multicolumn{1}{c}{\shortstack{$\max$\\$|\Delta\hat p|$}} & \multicolumn{1}{c}{$\Delta L$} & \multicolumn{1}{c}{$\Delta Y$} & & \multicolumn{1}{c}{$\Delta\omega$} & \multicolumn{1}{c}{\shortstack{$\max$\\$|\Delta\hat p|$}} & \multicolumn{1}{c}{$\Delta L$} & \multicolumn{1}{c}{$\Delta Y$} \\
\midrule
\multicolumn{16}{l}{\textbf{No-reallocation}} \\
\rowcolor{tableguide}\quad $0$ (no extension) &  & $+0.62$ & $22.14$ & $0.00$ & $-0.01$ &  & \multicolumn{9}{>{\columncolor{tableguide}}c}{\expandafter\let\csname CT@row@color\endcsname\relax\renewcommand{\arraystretch}{1}\begin{tabular}{@{}>{\raggedleft\arraybackslash}p{\pdwA}!{\color{tableline}\vrule width 0.3pt}>{\raggedleft\arraybackslash}p{\pdwB}!{\color{tableline}\vrule width 0.3pt}>{\raggedleft\arraybackslash}p{\pdwC}!{\color{tableline}\vrule width 0.3pt}>{\raggedleft\arraybackslash}p{\pdwD}@{}}$+0.70$ & $27.90$ & $0.00$ & $-0.02$\end{tabular}} \\
\quad $+0.5$ &  & $+0.64$ & $26.66$ & $0.00$ & $-0.09$ &  & \multicolumn{9}{c}{\expandafter\let\csname CT@row@color\endcsname\relax\renewcommand{\arraystretch}{1}\begin{tabular}{@{}>{\raggedleft\arraybackslash}p{\pdwA}!{\color{tableline}\vrule width 0.3pt}>{\raggedleft\arraybackslash}p{\pdwB}!{\color{tableline}\vrule width 0.3pt}>{\raggedleft\arraybackslash}p{\pdwC}!{\color{tableline}\vrule width 0.3pt}>{\raggedleft\arraybackslash}p{\pdwD}@{}}$+0.74$ & $35.30$ & $0.00$ & $-0.10$\end{tabular}} \\
\quad $-0.5$ &  & $+0.61$ & $16.59$ & $0.00$ & $+0.11$ &  & \multicolumn{9}{c}{\expandafter\let\csname CT@row@color\endcsname\relax\renewcommand{\arraystretch}{1}\begin{tabular}{@{}>{\raggedleft\arraybackslash}p{\pdwA}!{\color{tableline}\vrule width 0.3pt}>{\raggedleft\arraybackslash}p{\pdwB}!{\color{tableline}\vrule width 0.3pt}>{\raggedleft\arraybackslash}p{\pdwC}!{\color{tableline}\vrule width 0.3pt}>{\raggedleft\arraybackslash}p{\pdwD}@{}}$+0.67$ & $19.71$ & $0.00$ & $+0.11$\end{tabular}} \\
\addlinespace[6pt]
\multicolumn{16}{l}{\textbf{Standard CGE}} \\
\rowcolor{tableguide}\quad $0$ (no extension) &  & $0.00$ & $0.00$ & $0.00$ & $0.00$ &  & \multicolumn{9}{>{\columncolor{tableguide}}c}{\expandafter\let\csname CT@row@color\endcsname\relax\renewcommand{\arraystretch}{1}\begin{tabular}{@{}>{\raggedleft\arraybackslash}p{\pdwA}!{\color{tableline}\vrule width 0.3pt}>{\raggedleft\arraybackslash}p{\pdwB}!{\color{tableline}\vrule width 0.3pt}>{\raggedleft\arraybackslash}p{\pdwC}!{\color{tableline}\vrule width 0.3pt}>{\raggedleft\arraybackslash}p{\pdwD}@{}}$0.00$ & $0.00$ & $0.00$ & $0.00$\end{tabular}} \\
\quad $+0.5$ &  & $-0.12$ & $3.84$ & $0.00$ & $-0.17$ &  & \multicolumn{9}{c}{\expandafter\let\csname CT@row@color\endcsname\relax\renewcommand{\arraystretch}{1}\begin{tabular}{@{}>{\raggedleft\arraybackslash}p{\pdwA}!{\color{tableline}\vrule width 0.3pt}>{\raggedleft\arraybackslash}p{\pdwB}!{\color{tableline}\vrule width 0.3pt}>{\raggedleft\arraybackslash}p{\pdwC}!{\color{tableline}\vrule width 0.3pt}>{\raggedleft\arraybackslash}p{\pdwD}@{}}$-0.13$ & $4.05$ & $0.00$ & $-0.18$\end{tabular}} \\
\quad $-0.5$ &  & $+0.17$ & $4.33$ & $0.00$ & $+0.22$ &  & \multicolumn{9}{c}{\expandafter\let\csname CT@row@color\endcsname\relax\renewcommand{\arraystretch}{1}\begin{tabular}{@{}>{\raggedleft\arraybackslash}p{\pdwA}!{\color{tableline}\vrule width 0.3pt}>{\raggedleft\arraybackslash}p{\pdwB}!{\color{tableline}\vrule width 0.3pt}>{\raggedleft\arraybackslash}p{\pdwC}!{\color{tableline}\vrule width 0.3pt}>{\raggedleft\arraybackslash}p{\pdwD}@{}}$+0.16$ & $4.23$ & $0.00$ & $+0.22$\end{tabular}} \\
\addlinespace[6pt]
\multicolumn{16}{l}{\textbf{Elastic labor supply ($\eta_s = 0.5$)}} \\
\rowcolor{tableguide}\quad $0$ (no extension) &  & $+0.27$ & $9.58$ & $+0.12$ & $+0.12$ &  & \multicolumn{9}{>{\columncolor{tableguide}}c}{\expandafter\let\csname CT@row@color\endcsname\relax\renewcommand{\arraystretch}{1}\begin{tabular}{@{}>{\raggedleft\arraybackslash}p{\pdwA}!{\color{tableline}\vrule width 0.3pt}>{\raggedleft\arraybackslash}p{\pdwB}!{\color{tableline}\vrule width 0.3pt}>{\raggedleft\arraybackslash}p{\pdwC}!{\color{tableline}\vrule width 0.3pt}>{\raggedleft\arraybackslash}p{\pdwD}@{}}$+0.29$ & $10.57$ & $+0.13$ & $+0.13$\end{tabular}} \\
\quad $+0.5$ &  & $+0.23$ & $13.96$ & $+0.09$ & $-0.04$ &  & \multicolumn{9}{c}{\expandafter\let\csname CT@row@color\endcsname\relax\renewcommand{\arraystretch}{1}\begin{tabular}{@{}>{\raggedleft\arraybackslash}p{\pdwA}!{\color{tableline}\vrule width 0.3pt}>{\raggedleft\arraybackslash}p{\pdwB}!{\color{tableline}\vrule width 0.3pt}>{\raggedleft\arraybackslash}p{\pdwC}!{\color{tableline}\vrule width 0.3pt}>{\raggedleft\arraybackslash}p{\pdwD}@{}}$+0.25$ & $16.07$ & $+0.10$ & $-0.04$\end{tabular}} \\
\quad $-0.5$ &  & $+0.35$ & $4.43$ & $+0.16$ & $+0.36$ &  & \multicolumn{9}{c}{\expandafter\let\csname CT@row@color\endcsname\relax\renewcommand{\arraystretch}{1}\begin{tabular}{@{}>{\raggedleft\arraybackslash}p{\pdwA}!{\color{tableline}\vrule width 0.3pt}>{\raggedleft\arraybackslash}p{\pdwB}!{\color{tableline}\vrule width 0.3pt}>{\raggedleft\arraybackslash}p{\pdwC}!{\color{tableline}\vrule width 0.3pt}>{\raggedleft\arraybackslash}p{\pdwD}@{}}$+0.36$ & $4.65$ & $+0.17$ & $+0.37$\end{tabular}} \\
\addlinespace[6pt]
\multicolumn{16}{l}{\textbf{Employment adjustment}} \\
\rowcolor{tableguide}\quad $0$ (no extension) &  & $0.00$ & $0.00$ & $-0.10$ & $-0.10$ &  & $0.00$ & $0.00$ & $+0.13$ & $+0.13$ &  & $0.00$ & $0.00$ & $+1.78$ & $+1.78$ \\
\quad $+0.5$ &  & $-0.13$ & $3.86$ & $+0.12$ & $-0.06$ &  & $-0.15$ & $4.08$ & $+0.37$ & $+0.16$ &  & $-0.30$ & $4.26$ & $+2.21$ & $+1.86$ \\
\quad $-0.5$ &  & $+0.13$ & $4.26$ & $-0.32$ & $-0.13$ &  & $+0.15$ & $4.22$ & $-0.12$ & $+0.09$ &  & $+0.31$ & $4.39$ & $+1.35$ & $+1.71$ \\
\addlinespace[6pt]
\multicolumn{16}{l}{\textbf{Leontief closure}} \\
\rowcolor{tableguide}\quad $0$ (no extension) &  & $0.00$ & $0.00$ & $-0.10$ & $-0.10$ &  & $0.00$ & $0.00$ & $+0.13$ & $+0.13$ &  & $0.00$ & $0.00$ & $+1.78$ & $+1.78$ \\
\quad $+0.5$ &  & $-0.14$ & $4.27$ & $+0.09$ & $-0.10$ &  & $-0.17$ & $4.38$ & $+0.34$ & $+0.12$ &  & $-0.33$ & $4.58$ & $+2.13$ & $+1.77$ \\
\quad $-0.5$ &  & $+0.12$ & $3.86$ & $-0.26$ & $-0.10$ &  & $+0.14$ & $3.95$ & $-0.06$ & $+0.13$ &  & $+0.28$ & $4.10$ & $+1.48$ & $+1.78$ \\
\bottomrule
\end{tabular}}
\caption{Integrating a Kaldor--Verdoorn term in the general closure matrix. Each block reports
the average real consumption wage ($\Delta\omega$, the employment-weighted money wage deflated by
the consumer price index, Section~\ref{subsec:optimization}), the largest sectoral deviation from the
consumption basket ($\max|\Delta\hat p| = \max_i |p_i/\mathcal{P}-1|$), employment ($\Delta L$) and real GDP
($\Delta Y$, the wage bill plus the surplus, as in Table~\ref{tab:results-matrix}), all in percent against
the baseline. Both price statistics are independent of the numeraire and comparable across all
five panels. The shaded $\delta = 0$ rows (no extension) reproduce
results from Table~\ref{tab:results-matrix}.
The incremental effect of the extension is the difference between a row and that
shaded baseline.}
\label{tab:pricing-kaldor}
\end{table}


\subsection{Markup pricing and distributional assumptions}\label{subsec:pricing-markup}

Alongside the sectoral productivity effects discussed above, increasing demand may affect prices through utilization-dependent markups. Effective markups are said to rise with utilization mechanically as profits rise, while capital costs stay constant. Relatedly, it is argued that increased demand will lead to increased pricing power -- an intuition that shows a long pedigree in economics \citep{Kalecki1954,Bils1987}. Taking this basic idea -- markups rising with capacity utilization -- as a vantage point and projecting it onto a sectoral level in line with the basic setup from Section~\ref{sec:kernel}, we arrive at the following formulation for prices:

\begin{equation}\label{eq:markup-extension}
  p_i \;=\; \mu_i\,c_i, \qquad \mu_i \;=\; 1 + \kappa\Bigl(\frac{y_i}{\lambda_i} - 1\Bigr),
\end{equation}
which is nested at $\kappa = 0$ so that the core framework from Section~\ref{sec:kernel} remains unchanged. The parameter $\kappa$ has a direct reading reflecting our theoretical rationale: at $\kappa = 0.5$ a $10\%$ expansion lifts the margin by five percent, i.e., to $\mu_i = 1.05$. The largest margin realized across the runs presented below is $1.11$ (coinciding with a sectoral output response of about $22\%$). 

Moreover, introducing demand-responsive markups changes the income side of our model economies: Section~\ref{sec:kernel} reads value added from the wage bill, which represents composite factor income in our framework. Hence, the markup value imposed here drives an additional wedge between costs and profits that comes on top of the operating profits already included implicitly in the wage bill. In other words, such a pricing equation departs from the hitherto used benchmark of zero profits and introduces demand-driven profits as an explicit second income stream, without explicitly modeling capital as a second factor of production. Formally, this feature is introduced in a minimal way: firms minimize production costs at given output, while selling at markup prices $p_i=\mu_i c_i$ (see Appendix~\ref{app:pricing-hooks} for more details). Profits thus enter as a pure surplus (akin to the classic definition of a rent), and an additional assumption is required for where this income goes.
Written out, that surplus is
\begin{equation}\label{eq:profit-bill}
  \Pi \;=\; \sum_i \Bigl(1 - \frac{1}{\mu_i}\Bigr) p_i y_i ,
\end{equation}
the sum over sectors of revenue minus the marginal-cost bill, zero whenever
$\kappa = 0$ and positive exactly where prices exceed marginal cost.

The flow of profit income is governed by the retention share $s_\pi$: firms keep $s_\pi$ of the surplus and distribute the remainder, which enters the household budget as untaxed income. At $s_\pi = 1$ the rent is retained in full and booked as a leak beside household saving, so the household side is exactly that of the model core --- although prices and quantities still carry the markup. The distributed part, in contrast, is saved like any other income at $s \approx 0.12$. Table~\ref{tab:pricing-markup} reports full retention at $\kappa = 0.25$ and $\kappa = 0.5$ and two distributive variants at $\kappa = 0.5$, $s_\pi = 0.5$ and $s_\pi = 0$.\footnote{The augmented household budget is $E = (1-\tau)\sum_i w_i L_i + F + (1-s_\pi)\Pi$: the distributed share of the surplus enters beside after-tax labor income and the external transfer, and the household saves the share $s$ of the total as before, so $s_\pi = 1$ returns the model core's budget exactly. The overall saving rate on profits is therefore $1-(1-s)(1-s_\pi)$, which equals $s$ at $s_\pi = 0$ and $0.56$ at $s_\pi = 0.5$. The retained part $s_\pi\Pi$ is booked as a leak in the external account (Appendix~\ref{app:pricing-accounting}).}

% =====================================================================
% Table 7 --- the markup extension by labor closure and financing (GENERATED)
% ---------------------------------------------------------------------
% Do NOT edit by hand. Regenerate with:  make_pricing_door.py
% OpenCGE tag v0.5.0 (ADR-0044); source run ids are listed below.
% Sources: matrix_5x3-<row>-<fin>[-etas05],
% pricing_markup-<row>-<fin>-<k025|k05>,
% pricing_profit-<row>-<fin>-k05-<spi05|spi0>.
% No-extension kappa = 0 reference row, then retained and distributed profit,
% kappa in {0.25, 0.5}, all five labor rows (ADR-0032). kappa = 0.25
% distributed: tableE.
% NOTE (JK, resolved): the shading now marks the kappa = 0 baseline row, which
% restores the convention of Tables 4 and 6 (the baseline row is shaded).
% COLLAPSED CELLS: as in Table 6, F2 = F3 is printed once in the three
% flexible-wage panels, asserted on the PRINTED strings.
% =====================================================================
\begin{table}[htbp]
\centering
\footnotesize
\definecolor{tableguide}{gray}{0.92}
\definecolor{tableline}{gray}{0.72}
\setlength{\tabcolsep}{3pt}
\renewcommand{\arraystretch}{1.2}
\resizebox{\textwidth}{!}{%
\ifdefined\pdwA\else\newlength{\pdwA}\fi
\settowidth{\pdwA}{$+0.62$}\settowidth{\dimen0}{$\Delta\omega$}\ifdim\dimen0>\pdwA\setlength{\pdwA}{\dimen0}\fi
\ifdefined\pdwB\else\newlength{\pdwB}\fi
\settowidth{\pdwB}{$22.14$}\settowidth{\dimen0}{\shortstack{$\max$\\$|\Delta\hat p|$}}\ifdim\dimen0>\pdwB\setlength{\pdwB}{\dimen0}\fi
\ifdefined\pdwC\else\newlength{\pdwC}\fi
\settowidth{\pdwC}{$+0.12$}\settowidth{\dimen0}{$\Delta L$}\ifdim\dimen0>\pdwC\setlength{\pdwC}{\dimen0}\fi
\ifdefined\pdwD\else\newlength{\pdwD}\fi
\settowidth{\pdwD}{$-0.01$}\settowidth{\dimen0}{$\Delta Y$}\ifdim\dimen0>\pdwD\setlength{\pdwD}{\dimen0}\fi
\ifdefined\pdwE\else\newlength{\pdwE}\fi
\settowidth{\pdwE}{$0.00$}\settowidth{\dimen0}{$\frac{\Delta\Pi}{Y}$}\ifdim\dimen0>\pdwE\setlength{\pdwE}{\dimen0}\fi
\begin{tabular}{l@{\hspace{3pt}}p{2pt}@{\hspace{3pt}}>{\raggedleft\arraybackslash}p{\pdwA}!{\color{tableline}\vrule width 0.3pt}>{\raggedleft\arraybackslash}p{\pdwB}!{\color{tableline}\vrule width 0.3pt}>{\raggedleft\arraybackslash}p{\pdwC}!{\color{tableline}\vrule width 0.3pt}>{\raggedleft\arraybackslash}p{\pdwD}!{\color{tableline}\vrule width 0.3pt}>{\raggedleft\arraybackslash}p{\pdwE}@{\hspace{3pt}}p{2pt}@{\hspace{3pt}}>{\raggedleft\arraybackslash}p{\pdwA}!{\color{tableline}\vrule width 0.3pt}>{\raggedleft\arraybackslash}p{\pdwB}!{\color{tableline}\vrule width 0.3pt}>{\raggedleft\arraybackslash}p{\pdwC}!{\color{tableline}\vrule width 0.3pt}>{\raggedleft\arraybackslash}p{\pdwD}!{\color{tableline}\vrule width 0.3pt}>{\raggedleft\arraybackslash}p{\pdwE}@{\hspace{3pt}}p{2pt}@{\hspace{3pt}}>{\raggedleft\arraybackslash}p{\pdwA}!{\color{tableline}\vrule width 0.3pt}>{\raggedleft\arraybackslash}p{\pdwB}!{\color{tableline}\vrule width 0.3pt}>{\raggedleft\arraybackslash}p{\pdwC}!{\color{tableline}\vrule width 0.3pt}>{\raggedleft\arraybackslash}p{\pdwD}!{\color{tableline}\vrule width 0.3pt}>{\raggedleft\arraybackslash}p{\pdwE}}
\toprule
 & & \multicolumn{5}{c}{\textbf{F1}} & & \multicolumn{5}{c}{\textbf{F2}} & & \multicolumn{5}{c}{\textbf{F3}} \\
 & & \multicolumn{5}{c}{\scriptsize repurposing of consumption} & & \multicolumn{5}{c}{\scriptsize refinancing via taxes} & & \multicolumn{5}{c}{\scriptsize external finance} \\
\cmidrule(lr){3-7}\cmidrule(lr){9-13}\cmidrule(lr){15-19}
\multicolumn{1}{c}{Routing of profits} & & \multicolumn{1}{c}{$\Delta\omega$} & \multicolumn{1}{c}{\shortstack{$\max$\\$|\Delta\hat p|$}} & \multicolumn{1}{c}{$\Delta L$} & \multicolumn{1}{c}{$\Delta Y$} & \multicolumn{1}{c}{$\frac{\Delta\Pi}{Y}$} & & \multicolumn{1}{c}{$\Delta\omega$} & \multicolumn{1}{c}{\shortstack{$\max$\\$|\Delta\hat p|$}} & \multicolumn{1}{c}{$\Delta L$} & \multicolumn{1}{c}{$\Delta Y$} & \multicolumn{1}{c}{$\frac{\Delta\Pi}{Y}$} & & \multicolumn{1}{c}{$\Delta\omega$} & \multicolumn{1}{c}{\shortstack{$\max$\\$|\Delta\hat p|$}} & \multicolumn{1}{c}{$\Delta L$} & \multicolumn{1}{c}{$\Delta Y$} & \multicolumn{1}{c}{$\frac{\Delta\Pi}{Y}$} \\
\midrule
\multicolumn{19}{l}{\textbf{No-reallocation}} \\
\rowcolor{tableguide}\quad $\kappa = 0$ (no extension) &  & $+0.62$ & $22.14$ & $0.00$ & $-0.01$ & $0.00$ &  & \multicolumn{11}{>{\columncolor{tableguide}}c}{\expandafter\let\csname CT@row@color\endcsname\relax\renewcommand{\arraystretch}{1}\begin{tabular}{@{}>{\raggedleft\arraybackslash}p{\pdwA}!{\color{tableline}\vrule width 0.3pt}>{\raggedleft\arraybackslash}p{\pdwB}!{\color{tableline}\vrule width 0.3pt}>{\raggedleft\arraybackslash}p{\pdwC}!{\color{tableline}\vrule width 0.3pt}>{\raggedleft\arraybackslash}p{\pdwD}!{\color{tableline}\vrule width 0.3pt}>{\raggedleft\arraybackslash}p{\pdwE}@{}}$+0.70$ & $27.90$ & $0.00$ & $-0.02$ & $0.00$\end{tabular}} \\
\quad retained, $\kappa = 0.25$ &  & $+0.65$ & $24.43$ & $0.00$ & $-0.01$ & $0.04$ &  & \multicolumn{11}{c}{\expandafter\let\csname CT@row@color\endcsname\relax\renewcommand{\arraystretch}{1}\begin{tabular}{@{}>{\raggedleft\arraybackslash}p{\pdwA}!{\color{tableline}\vrule width 0.3pt}>{\raggedleft\arraybackslash}p{\pdwB}!{\color{tableline}\vrule width 0.3pt}>{\raggedleft\arraybackslash}p{\pdwC}!{\color{tableline}\vrule width 0.3pt}>{\raggedleft\arraybackslash}p{\pdwD}!{\color{tableline}\vrule width 0.3pt}>{\raggedleft\arraybackslash}p{\pdwE}@{}}$+0.73$ & $31.46$ & $0.00$ & $-0.02$ & $0.04$\end{tabular}} \\
\quad \begin{tabular}{@{}l@{}}retained, $\kappa = 0.5$\\distributed, $\kappa = 0.5$, $s_\pi = 0.5$\\distributed, $\kappa = 0.5$, $s_\pi = 0$\end{tabular} &  & $+0.68$ & $26.40$ & $0.00$ & $-0.02$ & $0.05$ &  & \multicolumn{11}{c}{\expandafter\let\csname CT@row@color\endcsname\relax\renewcommand{\arraystretch}{1}\begin{tabular}{@{}>{\raggedleft\arraybackslash}p{\pdwA}!{\color{tableline}\vrule width 0.3pt}>{\raggedleft\arraybackslash}p{\pdwB}!{\color{tableline}\vrule width 0.3pt}>{\raggedleft\arraybackslash}p{\pdwC}!{\color{tableline}\vrule width 0.3pt}>{\raggedleft\arraybackslash}p{\pdwD}!{\color{tableline}\vrule width 0.3pt}>{\raggedleft\arraybackslash}p{\pdwE}@{}}$+0.77$ & $34.63$ & $0.00$ & $-0.02$ & $0.06$\end{tabular}} \\
\addlinespace[6pt]
\multicolumn{19}{l}{\textbf{Standard CGE}} \\
\rowcolor{tableguide}\quad $\kappa = 0$ (no extension) &  & $0.00$ & $0.00$ & $0.00$ & $0.00$ & $0.00$ &  & \multicolumn{11}{>{\columncolor{tableguide}}c}{\expandafter\let\csname CT@row@color\endcsname\relax\renewcommand{\arraystretch}{1}\begin{tabular}{@{}>{\raggedleft\arraybackslash}p{\pdwA}!{\color{tableline}\vrule width 0.3pt}>{\raggedleft\arraybackslash}p{\pdwB}!{\color{tableline}\vrule width 0.3pt}>{\raggedleft\arraybackslash}p{\pdwC}!{\color{tableline}\vrule width 0.3pt}>{\raggedleft\arraybackslash}p{\pdwD}!{\color{tableline}\vrule width 0.3pt}>{\raggedleft\arraybackslash}p{\pdwE}@{}}$0.00$ & $0.00$ & $0.00$ & $0.00$ & $0.00$\end{tabular}} \\
\quad retained, $\kappa = 0.25$ &  & $+0.09$ & $5.32$ & $0.00$ & $-0.01$ & $0.07$ &  & \multicolumn{11}{c}{\expandafter\let\csname CT@row@color\endcsname\relax\renewcommand{\arraystretch}{1}\begin{tabular}{@{}>{\raggedleft\arraybackslash}p{\pdwA}!{\color{tableline}\vrule width 0.3pt}>{\raggedleft\arraybackslash}p{\pdwB}!{\color{tableline}\vrule width 0.3pt}>{\raggedleft\arraybackslash}p{\pdwC}!{\color{tableline}\vrule width 0.3pt}>{\raggedleft\arraybackslash}p{\pdwD}!{\color{tableline}\vrule width 0.3pt}>{\raggedleft\arraybackslash}p{\pdwE}@{}}$+0.09$ & $5.68$ & $0.00$ & $-0.01$ & $0.07$\end{tabular}} \\
\quad \begin{tabular}{@{}l@{}}retained, $\kappa = 0.5$\\distributed, $\kappa = 0.5$, $s_\pi = 0.5$\\distributed, $\kappa = 0.5$, $s_\pi = 0$\end{tabular} &  & $+0.17$ & $9.49$ & $0.00$ & $-0.01$ & $0.10$ &  & \multicolumn{11}{c}{\expandafter\let\csname CT@row@color\endcsname\relax\renewcommand{\arraystretch}{1}\begin{tabular}{@{}>{\raggedleft\arraybackslash}p{\pdwA}!{\color{tableline}\vrule width 0.3pt}>{\raggedleft\arraybackslash}p{\pdwB}!{\color{tableline}\vrule width 0.3pt}>{\raggedleft\arraybackslash}p{\pdwC}!{\color{tableline}\vrule width 0.3pt}>{\raggedleft\arraybackslash}p{\pdwD}!{\color{tableline}\vrule width 0.3pt}>{\raggedleft\arraybackslash}p{\pdwE}@{}}$+0.18$ & $10.50$ & $0.00$ & $-0.01$ & $0.11$\end{tabular}} \\
\addlinespace[6pt]
\multicolumn{19}{l}{\textbf{Elastic labor supply ($\eta_s = 0.5$)}} \\
\rowcolor{tableguide}\quad $\kappa = 0$ (no extension) &  & $+0.27$ & $9.58$ & $+0.12$ & $+0.12$ & $0.00$ &  & \multicolumn{11}{>{\columncolor{tableguide}}c}{\expandafter\let\csname CT@row@color\endcsname\relax\renewcommand{\arraystretch}{1}\begin{tabular}{@{}>{\raggedleft\arraybackslash}p{\pdwA}!{\color{tableline}\vrule width 0.3pt}>{\raggedleft\arraybackslash}p{\pdwB}!{\color{tableline}\vrule width 0.3pt}>{\raggedleft\arraybackslash}p{\pdwC}!{\color{tableline}\vrule width 0.3pt}>{\raggedleft\arraybackslash}p{\pdwD}!{\color{tableline}\vrule width 0.3pt}>{\raggedleft\arraybackslash}p{\pdwE}@{}}$+0.29$ & $10.57$ & $+0.13$ & $+0.13$ & $0.00$\end{tabular}} \\
\quad retained, $\kappa = 0.25$ &  & $+0.26$ & $13.43$ & $+0.12$ & $+0.11$ & $0.11$ &  & \multicolumn{11}{c}{\expandafter\let\csname CT@row@color\endcsname\relax\renewcommand{\arraystretch}{1}\begin{tabular}{@{}>{\raggedleft\arraybackslash}p{\pdwA}!{\color{tableline}\vrule width 0.3pt}>{\raggedleft\arraybackslash}p{\pdwB}!{\color{tableline}\vrule width 0.3pt}>{\raggedleft\arraybackslash}p{\pdwC}!{\color{tableline}\vrule width 0.3pt}>{\raggedleft\arraybackslash}p{\pdwD}!{\color{tableline}\vrule width 0.3pt}>{\raggedleft\arraybackslash}p{\pdwE}@{}}$+0.28$ & $15.35$ & $+0.12$ & $+0.12$ & $0.12$\end{tabular}} \\
\quad \begin{tabular}{@{}l@{}}retained, $\kappa = 0.5$\\distributed, $\kappa = 0.5$, $s_\pi = 0.5$\\distributed, $\kappa = 0.5$, $s_\pi = 0$\end{tabular} &  & $+0.26$ & $16.56$ & $+0.12$ & $+0.11$ & $0.18$ &  & \multicolumn{11}{c}{\expandafter\let\csname CT@row@color\endcsname\relax\renewcommand{\arraystretch}{1}\begin{tabular}{@{}>{\raggedleft\arraybackslash}p{\pdwA}!{\color{tableline}\vrule width 0.3pt}>{\raggedleft\arraybackslash}p{\pdwB}!{\color{tableline}\vrule width 0.3pt}>{\raggedleft\arraybackslash}p{\pdwC}!{\color{tableline}\vrule width 0.3pt}>{\raggedleft\arraybackslash}p{\pdwD}!{\color{tableline}\vrule width 0.3pt}>{\raggedleft\arraybackslash}p{\pdwE}@{}}$+0.29$ & $19.49$ & $+0.13$ & $+0.12$ & $0.20$\end{tabular}} \\
\addlinespace[6pt]
\multicolumn{19}{l}{\textbf{Employment adjustment}} \\
\rowcolor{tableguide}\quad $\kappa = 0$ (no extension) &  & $0.00$ & $0.00$ & $-0.10$ & $-0.10$ & $0.00$ &  & $0.00$ & $0.00$ & $+0.13$ & $+0.13$ & $0.00$ &  & $0.00$ & $0.00$ & $+1.78$ & $+1.78$ & $0.00$ \\
\quad retained, $\kappa = 0.25$ &  & $+0.08$ & $5.32$ & $0.00$ & $0.00$ & $0.07$ &  & $+0.01$ & $5.65$ & $+0.16$ & $+0.15$ & $0.14$ &  & $-0.58$ & $5.51$ & $+1.28$ & $+1.24$ & $0.62$ \\
\quad retained, $\kappa = 0.5$ &  & $+0.12$ & $9.49$ & $+0.06$ & $+0.05$ & $0.15$ &  & $+0.01$ & $10.45$ & $+0.18$ & $+0.17$ & $0.25$ &  & $-0.83$ & $10.21$ & $+1.08$ & $+1.02$ & $0.95$ \\
\quad distributed, $\kappa = 0.5$, $s_\pi = 0.5$ &  & $+0.06$ & $9.49$ & $+0.12$ & $+0.11$ & $0.20$ &  & $-0.09$ & $10.42$ & $+0.29$ & $+0.27$ & $0.34$ &  & $-1.21$ & $10.11$ & $+1.49$ & $+1.42$ & $1.27$ \\
\quad distributed, $\kappa = 0.5$, $s_\pi = 0$ &  & $-0.07$ & $9.48$ & $+0.25$ & $+0.23$ & $0.30$ &  & $-0.30$ & $10.36$ & $+0.52$ & $+0.48$ & $0.51$ &  & $-1.99$ & $9.89$ & $+2.35$ & $+2.24$ & $1.93$ \\
\addlinespace[6pt]
\multicolumn{19}{l}{\textbf{Leontief closure}} \\
\rowcolor{tableguide}\quad $\kappa = 0$ (no extension) &  & $0.00$ & $0.00$ & $-0.10$ & $-0.10$ & $0.00$ &  & $0.00$ & $0.00$ & $+0.13$ & $+0.13$ & $0.00$ &  & $0.00$ & $0.00$ & $+1.78$ & $+1.78$ & $0.00$ \\
\quad retained, $\kappa = 0.25$ &  & $+0.10$ & $6.16$ & $-0.03$ & $-0.03$ & $0.08$ &  & $+0.02$ & $6.26$ & $+0.11$ & $+0.11$ & $0.15$ &  & $-0.59$ & $6.12$ & $+1.12$ & $+1.12$ & $0.65$ \\
\quad retained, $\kappa = 0.5$ &  & $+0.16$ & $12.44$ & $0.00$ & $0.00$ & $0.20$ &  & $+0.04$ & $12.64$ & $+0.10$ & $+0.10$ & $0.30$ &  & $-0.84$ & $12.41$ & $+0.84$ & $+0.84$ & $1.02$ \\
\quad distributed, $\kappa = 0.5$, $s_\pi = 0.5$ &  & $+0.07$ & $12.44$ & $+0.07$ & $+0.07$ & $0.27$ &  & $-0.09$ & $12.61$ & $+0.21$ & $+0.21$ & $0.41$ &  & $-1.28$ & $12.29$ & $+1.21$ & $+1.20$ & $1.38$ \\
\quad distributed, $\kappa = 0.5$, $s_\pi = 0$ &  & $-0.11$ & $12.44$ & $+0.22$ & $+0.22$ & $0.42$ &  & $-0.36$ & $12.53$ & $+0.44$ & $+0.44$ & $0.63$ &  & $-2.19$ & $12.05$ & $+1.98$ & $+1.97$ & $2.13$ \\
\bottomrule
\end{tabular}}
\caption{Introducing markup pricing in the context of a demand-only program for all closure scenarios.
Columns, real-wage and dispersion conventions as in Table~\ref{tab:pricing-kaldor}, including real
GDP ($\Delta Y$, the wage bill plus the surplus); the final
column, $\Delta\Pi/Y$, is the surplus the extension creates as a share of value added, in
percent (the model core prices at marginal
cost, so the baseline surplus is zero). As the shading indicates, the $\kappa = 0$ rows
reproduce the real-wage and employment cells of Table~\ref{tab:results-matrix} together with
the $\Delta Y$ column of Table~\ref{tab:pricing-kaldor}. The incremental effect of
the extension is the difference between a row and that shaded baseline.}
\label{tab:pricing-markup}
\end{table}


A first observation is that the direction of the real-wage response to markup pricing depends on the closure. In the three flexible-wage panels markups leave the average real wage unchanged or raise it slightly --- by $0.18$ percentage points in the standard CGE model at $\kappa = 0.5$ --- because the markup rule is symmetric: sectors operating below their baseline output price below cost, and these are mainly the consumer-facing sectors that the program crowds out, so the consumption basket becomes cheaper relative to the wage while the program's suppliers charge above cost. In the wage-pinned panels the money wage cannot adjust, and with profits retained the average real wage rises slightly under repurposing (by up to $0.16\%$), barely moves under tax finance and falls by about $0.8\%$ under external finance. These changes remain below one percent in the retained scenarios and have to be contrasted with $\max_i |p_i/\mathcal{P} - 1|$ --- the two leftmost columns of each panel --- where strong price changes are associated with sectors affected by the demand shock (see also Figure~\ref{fig:sectoral-comparison-scatter}). At $\kappa = 0.5$ the markup's largest sectoral deviation lies between $9.5\%$ and $35\%$ for the flexible wage models and between $9\%$ and $13\%$ for the wage-pinned models. However, as the former already exhibit wide sectoral price dispersion without markups --- the $\kappa = 0$ row alone moves prices by up to $27.9\%$ in the no-reallocation panel --- the relative impact of markup pricing on the largest sectoral deviation is larger for the latter group, which start from zero dispersion. Nonetheless, the introduction of demand-driven markups still raises the largest sectoral price change by up to 10.5 percentage points for the standard CGE model. In the no-reallocation panel the comparison is starkest: the average real wage already rises by $0.70\%$ without markups, and the markup at $\kappa = 0.5$ adds only $0.07$ percentage points with a profit bill of $0.06\%$ of value added, so almost the whole movement there reflects the dispersion of sectoral wages rather than the markup.

A second structural feature of these results is that, within a given financing column, the outcomes for the three flexible-wage models depend only on the intensity of markups ($\kappa$), not on how they are used, because the net external transfer $F$ adjusts and falls when the share of distributed profits rises, and vice versa. Distributing profits, by contrast, raises GDP and employment relative to retention in every financing column of the employment-adjustment and Leontief panels, most so under external finance. Under external financing (F3) at $\kappa = 0.5$, distributing rather than retaining profits deepens the fall in the average real wage from $-0.83\%$ to $-1.99\%$ under employment adjustment and from $-0.84\%$ to $-2.19\%$ under Leontief (Table~\ref{tab:pricing-markup}): with the money wage pinned, consumer prices rise by $2.04\%$ and $2.24\%$, and households lose about two percent of the purchasing power of their wage. Despite this decline in the real wage, employment and GDP rise with profit distribution.

The reason for the stronger employment effects observed in scenarios with distributed profits and fixed wages is simply that extra profit is spent rather than hoarded, so that effective multipliers become correspondingly larger. As with the Kaldor--Verdoorn effects, the coincidence between the employment-adjustment and Leontief columns no longer holds once prices move: in this context, the less flexible Leontief closures show different price changes and smaller growth of employment and output --- also relative to their own baselines.

Real GDP, which now includes the surplus $\Pi$ alongside the wage bill, falls from $1.78\%$ without the extension to $1.02\%$ under employment adjustment and $0.84\%$ under Leontief at $\kappa = 0.5$, and rises to $2.24\%$ and $1.97\%$ once the surplus is distributed. One caveat bounds this reading: the markup is symmetric about the baseline, so a contracting sector prices at a discount ($\mu_i < 1$).


